\chapter{Complexity — Single Tape}
%%% generated by the blueprint pipeline from TCSlib.Complexity.TuringMachine.Robustness.SingleTape
\section{Overview}

This module proves the in-model analogue of [AB09, Claim 1.6]: a Turing machine with $k$
work tapes over a finite alphabet $\Gamma$ (and separate read-only input and write-only
output tapes) is simulated by a machine with a \emph{single} work tape over an enlarged
finite alphabet, with quadratic slowdown $c \cdot (T(n)+1)^2$. The single tape stores the
$k$ source tapes interleaved, with head positions marked in the cells; each source step
is simulated by a forward reading sweep and a backward writing sweep across the visited
zone. Combined with alphabet reduction ([AB09, Claim 1.5]) this yields a one-work-tape
machine over the binary alphabet with the same quadratic slowdown.

\section{Declarations}

\begin{definition}[Adding one unused work tape]
\lean{Turing.FinTM.unusedTapeTM}
\label{Turing.FinTM.unusedTapeTM}
\leanok
\uses{Turing.FinTM}
Let $M$ be a Turing machine over an alphabet $\Gamma$, with a read-only input tape, a
write-only output tape, and no work tapes. This constructs a machine with the same
states and exactly one work tape: each transition performs the input-head movement,
output emission, and state change of the corresponding transition of $M$, while writing
nothing on the added work tape and keeping its head stationary.
\end{definition}

\begin{definition}[Configuration embedding for the unused tape]
\lean{Turing.FinTM.unusedTapeCfg}
\label{Turing.FinTM.unusedTapeCfg}
\leanok
\uses{Turing.Cfg, Turing.FinTM}
For a Turing machine $M$ over an alphabet $\Gamma$ and a configuration $c$ of $M$ on an
input $x$, this produces the configuration of a one-work-tape machine with the same
state, input-head position, and output as $c$, whose single work tape is entirely blank
with its head at the origin.
\end{definition}

\begin{lemma}[Unused tape commutes with one step]
\lean{Turing.FinTM.unusedTape_step}
\label{Turing.FinTM.unusedTape_step}
\leanok
\uses{Turing.Cfg, Turing.Cfg.workTapeSymbols, Turing.FinTM, Turing.FinTM.SweepCell, Turing.FinTM.unusedTapeCfg, Turing.FinTM.unusedTapeTM, Turing.MultiTapeTM.step}
\difficulty{4}
Let $M$ be a Turing machine over an alphabet $\Gamma$ with no work tapes, and let $M'$
be the machine obtained from $M$ by adding one unused work tape. For every configuration
$c$ of $M$ on an input $x$, one transition of $M'$ applied to the embedding of $c$ (same
state, input position, and output; blank added tape with head at the origin) equals the
embedding of the configuration obtained by one transition of $M$ from $c$; this also
holds when $c$ is halted.
\end{lemma}

\begin{lemma}[Unused tape preserves computation]
\lean{Turing.FinTM.unusedTape_computes}
\label{Turing.FinTM.unusedTape_computes}
\leanok
\uses{Turing.FinTM, Turing.FinTM.ComputesFunInTime, Turing.FinTM.SweepCell, Turing.FinTM.computesInTime_iff, Turing.FinTM.unusedTapeCfg, Turing.FinTM.unusedTapeTM, Turing.FinTM.unusedTape_step, Turing.MultiTapeTM.initCfg, Turing.MultiTapeTM.runFrom_comm_of_step}
\difficulty{3}
Let $M$ be a Turing machine over an alphabet $\Gamma$ with no work tapes that computes a
function $f$ on words over $\Gamma$ within time $T$, i.e.\ on every input $x$ the run of
$M$ has halted after $T(\abs{x})$ steps with output $f(x)$. Then the machine obtained
from $M$ by adding one unused work tape also computes $f$ within time $T$.
\end{lemma}

\begin{definition}[Cell of the interleaved tape encoding]
\lean{Turing.FinTM.SweepCell}
\label{Turing.FinTM.SweepCell}
\leanok
For an alphabet $\Gamma$ and $k \in \mathbb{N}$, a sweep cell is a tuple consisting of a
tape index $i \in \{0, \dots, k-1\}$, an optional payload symbol from $\Gamma$ (absence
meaning blank), a Boolean head flag, and a second Boolean flag in which the forward
sweep records the head flag of the corresponding cell of the preceding block.
\end{definition}

\begin{definition}[Forward-sweep visit rule]
\lean{Turing.FinTM.readVisit}
\label{Turing.FinTM.readVisit}
\leanok
\uses{Turing.FinTM.SweepCell, Turing.FinTM.indexedVisit}
The visit rule of the forward sweep. Its control state is a table assigning to each tape
index an optional symbol and a Boolean flag. On visiting a sweep cell for tape $i$, the
table entry for $i$ becomes the cell's payload if the cell's head flag is set (otherwise
the recorded symbol is unchanged), paired with that head flag; the cell keeps its
payload and head flag and stores the Boolean previously held in the table, which is the
head flag of tape $i$ in the preceding block.
\end{definition}

\begin{definition}[Backward-sweep visit rule]
\lean{Turing.FinTM.writeVisit}
\label{Turing.FinTM.writeVisit}
\leanok
\uses{Turing.Action, Turing.FinTM.SweepCell, Turing.FinTM.indexedVisit}
The visit rule of the backward sweep, parametrised by an action $\alpha$ of the
simulated $k$-work-tape machine. Its control state is a table of Booleans, one per tape
(the head flags seen in the block to the right). On visiting a sweep cell for tape $i$:
the new table records the cell's head flag; the payload is overwritten by the symbol
$\alpha$ writes on tape $i$ when this cell carries the head (and kept otherwise); the
new head flag is taken from the right-neighbour table, from the current cell, or from
the stored left-neighbour flag according as $\alpha$ moves head $i$ left, keeps it in
place, or moves it right; and the stored left-neighbour flag is cleared.
\end{definition}

\begin{definition}[Head-position flag]
\lean{Turing.FinTM.headAt}
\label{Turing.FinTM.headAt}
\leanok
\uses{Turing.Cfg}
For a configuration $c$ of a machine with $k$ work tapes, a tape index $i$, and a
coordinate $j \in \mathbb{Z}$, the Boolean that is true exactly when the head of work
tape $i$ in $c$ is at coordinate $j$.
\end{definition}

\begin{definition}[Interleaved block at a coordinate]
\lean{Turing.FinTM.tapeRow}
\label{Turing.FinTM.tapeRow}
\leanok
\uses{Turing.Cfg, Turing.FinTM.SweepCell, Turing.FinTM.headAt}
The block of $k$ sweep cells encoding coordinate $j \in \mathbb{Z}$ of a configuration
$c$ of a $k$-work-tape machine: the cell for tape $i$ carries the symbol of tape $i$ at
coordinate $j$ and the flag indicating whether head $i$ is at $j$. When a Boolean
parameter is set, the cell additionally records whether head $i$ is at $j-1$; otherwise
that field is false.
\end{definition}

\begin{definition}[Forward-sweep control table]
\lean{Turing.FinTM.readState}
\label{Turing.FinTM.readState}
\leanok
\uses{Turing.Cfg, Turing.Cfg.workTapeSymbols, Turing.FinTM.headAt}
The control table that the forward sweep holds immediately before reading the block at
coordinate $j$ in a configuration $c$: for each tape $i$ it contains the symbol
currently scanned by head $i$ if that head lies at a coordinate strictly less than $j$
(and nothing otherwise), together with the flag indicating whether head $i$ is at
coordinate $j-1$.
\end{definition}

\begin{lemma}[Forward sweep over one block]
\lean{Turing.FinTM.read_row}
\label{Turing.FinTM.read_row}
\leanok
\uses{Turing.Cfg, Turing.Cfg.workTapeSymbols, Turing.FinTM.headAt, Turing.FinTM.indexedFold_block, Turing.FinTM.readState, Turing.FinTM.readVisit, Turing.FinTM.sweepFold, Turing.FinTM.tapeRow}
\difficulty{4}
Let $c$ be a configuration of a $k$-work-tape machine and $j \in \mathbb{Z}$.
Transducing the block encoding coordinate $j$ (with unrecorded left-neighbour flags) by
the forward-sweep visit rule, starting from the control table associated with $j$,
yields exactly the control table associated with $j+1$ together with the same block with
left-neighbour head flags recorded.
\end{lemma}

\begin{lemma}[Backward sweep over one block]
\lean{Turing.FinTM.write_row}
\label{Turing.FinTM.write_row}
\leanok
\uses{Turing.Action, Turing.Action.apply, Turing.Cfg, Turing.FinTM.headAt, Turing.FinTM.indexedFold_block_reverse, Turing.FinTM.sweepFold, Turing.FinTM.tapeRow, Turing.FinTM.writeVisit}
\difficulty{5}
Let $c$ be a configuration of a $k$-work-tape machine, $\alpha$ an action of that
machine, and $j \in \mathbb{Z}$. Folding the backward-sweep visit rule for $\alpha$ over
the reversed block encoding coordinate $j$ (with recorded left-neighbour flags),
starting from the table of head flags at coordinate $j+1$, yields the table of head
flags at coordinate $j$ together with the reversed block encoding coordinate $j$ of the
configuration obtained by applying $\alpha$ to $c$.
\end{lemma}

\begin{definition}[Zone of consecutive blocks]
\lean{Turing.FinTM.tapeZone}
\label{Turing.FinTM.tapeZone}
\leanok
Given a family of blocks indexed by integer coordinates and a start coordinate $j$, the
zone of size $n$ is the concatenation of the $n$ consecutive blocks at coordinates
$j, j+1, \dots, j+n-1$, in ascending order.
\end{definition}

\begin{lemma}[Forward sweep over a zone]
\lean{Turing.FinTM.read_zone}
\label{Turing.FinTM.read_zone}
\leanok
\uses{Turing.Cfg, Turing.FinTM.readState, Turing.FinTM.readVisit, Turing.FinTM.read_row, Turing.FinTM.sweepFold, Turing.FinTM.sweepFold_append, Turing.FinTM.tapeRow, Turing.FinTM.tapeZone}
\difficulty{4}
Let $c$ be a configuration of a $k$-work-tape machine, $j \in \mathbb{Z}$, and
$n \in \mathbb{N}$. The forward sweep over the concatenation of the $n$ unread blocks at
coordinates $j, \dots, j+n-1$, started from the control table at $j$, ends in the
control table at $j+n$ and produces the corresponding blocks with left-neighbour head
flags recorded.
\end{lemma}

\begin{lemma}[Backward sweep over a zone]
\lean{Turing.FinTM.write_zone}
\label{Turing.FinTM.write_zone}
\leanok
\uses{Turing.Action, Turing.Action.apply, Turing.Cfg, Turing.FinTM.SweepAlphabet, Turing.FinTM.headAt, Turing.FinTM.sweepFold, Turing.FinTM.sweepFold_append, Turing.FinTM.tapeRow, Turing.FinTM.tapeZone, Turing.FinTM.writeVisit, Turing.FinTM.write_row}
\difficulty{4}
Let $c$ be a configuration of a $k$-work-tape machine, $\alpha$ an action of that
machine, $j \in \mathbb{Z}$, and $n \in \mathbb{N}$. The backward sweep for $\alpha$
over the reversed concatenation of the $n$ read-marked blocks at coordinates
$j, \dots, j+n-1$, started from the table of head flags at coordinate $j+n$, ends in the
table of head flags at $j$ and produces the reversed unread blocks of the configuration
obtained by applying $\alpha$ to $c$.
\end{lemma}

\begin{definition}[Enlarged sweep alphabet]
\lean{Turing.FinTM.SweepAlphabet}
\label{Turing.FinTM.SweepAlphabet}
\leanok
\uses{Turing.FinTM.SweepCell}
The enlarged alphabet of the single-tape simulator over a source alphabet $\Gamma$ with
$k$ tapes: a symbol is either a data symbol from $\Gamma$ (used unchanged for input and
output), or an internal sweep cell, or a distinguished boundary marker. Formally it is
the disjoint union of $\Gamma$ with the set of sweep cells extended by one extra
element.
\end{definition}

\begin{definition}[Embedding of the source alphabet]
\lean{Turing.FinTM.sweepEmbed}
\label{Turing.FinTM.sweepEmbed}
\leanok
\uses{Turing.FinTM.SweepAlphabet}
The injection of the source alphabet $\Gamma$ into the enlarged sweep alphabet, sending
each symbol to the corresponding data symbol.
\end{definition}

\begin{definition}[Boundary marker]
\lean{Turing.FinTM.sweepBoundary}
\label{Turing.FinTM.sweepBoundary}
\leanok
\uses{Turing.FinTM.SweepAlphabet}
The distinguished boundary marker of the enlarged sweep alphabet: a non-blank symbol
distinct from every data symbol and from every internal cell, including cells with blank
payload.
\end{definition}

\begin{definition}[Internal cell symbol]
\lean{Turing.FinTM.sweepSymbol}
\label{Turing.FinTM.sweepSymbol}
\leanok
\uses{Turing.FinTM.SweepAlphabet, Turing.FinTM.SweepCell}
The map tagging a sweep cell as an internal symbol of the enlarged sweep alphabet.
\end{definition}

\begin{definition}[Decoding a scanned input symbol]
\lean{Turing.FinTM.sweepInput}
\label{Turing.FinTM.sweepInput}
\leanok
\uses{Turing.FinTM.SweepAlphabet}
The partial decoding of an optionally scanned enlarged symbol: a data symbol decodes to
the underlying symbol of $\Gamma$, while everything else (a blank, an internal cell, or
a boundary marker) decodes to nothing.
\end{definition}

\begin{definition}[Simulator control states]
\lean{Turing.FinTM.SweepState}
\label{Turing.FinTM.SweepState}
\leanok
The finite state space of the single-tape simulator of a $k$-work-tape machine with
state set $S$ over $\Gamma$. Its phases are: initialization, carrying a counter in
$\{0, \dots, k\}$; a return-to-origin walk; left extension of the zone, carrying a
source state and a counter; the forward reading sweep, carrying a source state and a
table assigning to each tape an optional symbol and a Boolean; right extension,
additionally carrying a counter; and the backward writing sweep, carrying a source
state, the scanned input symbol, the collected reads, and a table of Booleans. No
unbounded tape coordinate enters the state, so the state space is finite.
\end{definition}

\begin{definition}[Stationary phase-change action]
\lean{Turing.FinTM.sweepMove}
\label{Turing.FinTM.sweepMove}
\leanok
\uses{Turing.Action}
For a one-work-tape machine, the action that keeps the input head stationary, writes
nothing on the work tape (preserving the scanned cell) while moving its head in a given
direction $d$, emits no output, and makes a given optional state the next state.
\end{definition}

\begin{definition}[The single-tape sweep simulator]
\lean{Turing.FinTM.sweepTM}
\label{Turing.FinTM.sweepTM}
\leanok
\uses{Turing.FinTM, Turing.FinTM.SweepAlphabet, Turing.FinTM.SweepState, Turing.FinTM.readVisit, Turing.FinTM.sweepAct, Turing.FinTM.sweepBoundary, Turing.FinTM.sweepInput, Turing.FinTM.sweepMove, Turing.FinTM.sweepSymbol, Turing.FinTM.writeVisit}
The single-work-tape simulator of a machine $M$ with $k$ work tapes over a finite
alphabet $\Gamma$: a machine over the enlarged sweep alphabet, with one work tape, whose
finite control realises the phases of the simulation. Initialization writes the marked
origin block and its two boundary markers. Each simulated step extends the stored zone
by a blank block on the left, sweeps forward reading the marked payloads while recording
left-neighbour head flags in the cells, extends the zone on the right, performs the
source transition on the input head and output, and sweeps backward rewriting the cells,
carrying right-neighbour head flags in its control.
\end{definition}

\begin{definition}[Blank block]
\lean{Turing.FinTM.blankRow}
\label{Turing.FinTM.blankRow}
\leanok
\uses{Turing.FinTM.SweepCell}
The blank block of $k$ sweep cells: the cell for tape $i$ carries index $i$, no payload,
a head flag equal to a given mark, and a cleared left-neighbour flag. With the mark set
it is the initialized origin block; with the mark cleared it is a fresh guard block.
\end{definition}

\begin{lemma}[Blocks outside the support are blank]
\lean{Turing.FinTM.tapeRow_blank}
\label{Turing.FinTM.tapeRow_blank}
\leanok
\uses{Turing.Cfg, Turing.FinTM.blankRow, Turing.FinTM.headAt, Turing.FinTM.tapeRow}
\difficulty{2}
Let $c$ be a configuration of a $k$-work-tape machine and $j \in \mathbb{Z}$ a
coordinate such that no work-tape head of $c$ is at $j$ and every work tape of $c$ is
blank at $j$. Then the unread block encoding coordinate $j$ equals the blank unmarked
block.
\end{lemma}

\begin{lemma}[Block length]
\lean{Turing.FinTM.tapeRow_length}
\label{Turing.FinTM.tapeRow_length}
\leanok
\uses{Turing.Cfg, Turing.FinTM.tapeRow}
\difficulty{1}
Every block encoding a single coordinate of a configuration of a $k$-work-tape machine
is a list of length $k$, in both its unread and read-marked forms.
\end{lemma}

\begin{lemma}[Zone length]
\lean{Turing.FinTM.tapeZone_length}
\label{Turing.FinTM.tapeZone_length}
\leanok
\uses{Turing.Cfg, Turing.FinTM.tapeRow, Turing.FinTM.tapeRow_length, Turing.FinTM.tapeZone}
\difficulty{3}
A zone of $n$ consecutive blocks of a configuration of a $k$-work-tape machine has
length $n \cdot k$.
\end{lemma}

\begin{lemma}[Splitting a zone at a block boundary]
\lean{Turing.FinTM.tapeZone_append}
\label{Turing.FinTM.tapeZone_append}
\leanok
\uses{Turing.FinTM.tapeZone}
\difficulty{4}
For any family of blocks indexed by $\mathbb{Z}$ and any $j \in \mathbb{Z}$ and
$n, m \in \mathbb{N}$, the zone of $n+m$ blocks starting at $j$ is the concatenation of
the zone of $n$ blocks starting at $j$ with the zone of $m$ blocks starting at $j+n$.
\end{lemma}

\begin{definition}[Transport of input-head positions]
\lean{Turing.FinTM.sweepPos}
\label{Turing.FinTM.sweepPos}
\leanok
\uses{Turing.FinTM.SweepAlphabet}
For an input word $x$ over $\Gamma$, the transport of an input-head position on $x$ (one
of $\abs{x}+2$ positions, including the two endmarkers) to the numerically identical
position on the symbol-wise encoded input, which has the same length.
\end{definition}

\begin{lemma}[Input movement commutes with transport]
\lean{Turing.FinTM.sweepPos_move}
\label{Turing.FinTM.sweepPos_move}
\leanok
\uses{Turing.FinTM.sweepPos, Turing.moveInputPos}
\difficulty{2}
Moving the input head one step in any direction commutes with the transport of head
positions from an input word to its symbol-wise encoding.
\end{lemma}

\begin{lemma}[Scanned input decodes correctly]
\lean{Turing.FinTM.sweepInput_read}
\label{Turing.FinTM.sweepInput_read}
\leanok
\uses{Turing.Cfg, Turing.Cfg.inputSymbol, Turing.FinTM.SweepAlphabet, Turing.FinTM.sweepInput, Turing.FinTM.sweepPos}
\difficulty{4}
Let $c$ be a configuration of a $k$-work-tape machine on input $x$, and let $d$ be a
configuration of any machine over the enlarged sweep alphabet on the symbol-wise encoded
input whose input head is at the transported position of $c$'s input head. Then decoding
the input symbol scanned by $d$ yields exactly the input symbol scanned by $c$.
\end{lemma}

\begin{definition}[Canonical between-steps configuration]
\lean{Turing.FinTM.sweepStart}
\label{Turing.FinTM.sweepStart}
\leanok
\uses{Turing.Cfg, Turing.FinTM, Turing.FinTM.SweepAlphabet, Turing.FinTM.SweepState, Turing.FinTM.sweepBoundary, Turing.FinTM.sweepPos, Turing.FinTM.sweepRevCfg, Turing.FinTM.sweepSymbol, Turing.FinTM.tapeRow, Turing.FinTM.tapeZone}
The canonical simulator configuration representing a source configuration $c$ between
simulated steps, parametrised by a start coordinate $j$, a block count $n$, and an
absolute cell $z$. The control is in the left-extension phase carrying the state of $c$
(halted when $c$ is halted); the input head is at the transported position of $c$'s
input head; the output holds the symbol-wise encoding of $c$'s output. The work-tape
head scans a boundary marker at cell $z$; the cells above it hold the encoded blocks for
coordinates $j, j+1, \dots, j+n-1$ in ascending order, closed by a second boundary
marker, and all remaining cells are blank.
\end{definition}

\begin{lemma}[Writing at the focus of a zipper tape]
\lean{Turing.FinTM.sweepTape_write}
\label{Turing.FinTM.sweepTape_write}
\leanok
\uses{Turing.FinTM.sweepTape, Turing.FinTM.sweepTape_read}
\difficulty{4}
Consider the two-way infinite partial tape assembled from a position $z \in \mathbb{Z}$,
a list $r$ of optional symbols occupying cells $z, z+1, \dots$, and a list $l$ occupying
cells $z-1, z-2, \dots$ (all unlisted cells blank). Overwriting the cell at $z$ with a
value $b$ yields the tape assembled from $z$, $l$, and the list whose first entry is $b$
followed by the tail of $r$.
\end{lemma}

\begin{lemma}[Turning from a forward into a backward scan]
\lean{Turing.FinTM.sweep_turn_left}
\label{Turing.FinTM.sweep_turn_left}
\leanok
\uses{Turing.Action, Turing.Action.apply, Turing.FinTM.sweepCfg, Turing.FinTM.sweepRevCfg, Turing.FinTM.sweepTape, Turing.FinTM.sweepTape_turn, Turing.FinTM.sweepTape_write, Turing.moveInputPos}
\difficulty{3}
Consider a one-work-tape configuration in forward-facing zipper form: work head at cell
$z$, a list $r$ occupying cells $z, z+1, \dots$, a list $l$ occupying cells
$z-1, z-2, \dots$, an input-head position, and an accumulated output. Applying an action
that writes a symbol $b$ at the current cell, moves the work head left, moves the input
head in a given direction, optionally emits one output symbol, and enters a given
optional state produces the backward-facing zipper at $z-1$: cells $z, z+1, \dots$ hold
$b$ followed by the tail of $r$, cells $z-1, z-2, \dots$ hold $l$, the input head has
moved in the given direction, and the output is extended by the emitted symbol.
\end{lemma}

\begin{lemma}[Turning from a backward into a forward scan]
\lean{Turing.FinTM.sweep_turn_right}
\label{Turing.FinTM.sweep_turn_right}
\leanok
\uses{Turing.Action.apply, Turing.FinTM.sweepAct, Turing.FinTM.sweepCfg, Turing.FinTM.sweepRevCfg, Turing.FinTM.sweepTape, Turing.FinTM.sweepTape_turn, Turing.FinTM.sweepTape_write, Turing.moveInputPos_zero}
\difficulty{4}
Consider a one-work-tape configuration in backward-facing zipper form: work head at cell
$z$, a list $r$ occupying cells $z, z-1, \dots$, and a list $l$ occupying cells
$z+1, z+2, \dots$. Applying the sweep action that writes a symbol $b$ at the current
cell, moves the work head right, leaves the input head and output unchanged, and enters
a state $q'$ yields the forward-facing zipper at $z+1$: cells $z, z-1, \dots$ hold $b$
followed by the tail of $r$, and cells $z+1, z+2, \dots$ hold $l$.
\end{lemma}

\begin{lemma}[Initialization writes the origin block]
\lean{Turing.FinTM.sweep_init_block}
\label{Turing.FinTM.sweep_init_block}
\leanok
\uses{Turing.Action.apply, Turing.FinTM, Turing.FinTM.SweepAlphabet, Turing.FinTM.blankRow, Turing.FinTM.sweepCfg, Turing.FinTM.sweepCfg_right_any, Turing.FinTM.sweepSymbol, Turing.FinTM.sweepTM, Turing.FinTM.sweep_generate, Turing.MultiTapeTM.runFrom, Turing.MultiTapeTM.step}
\difficulty{4}
For the single-tape simulator of a machine $M$ with $k$ work tapes: from any
forward-facing zipper configuration whose control is in the initialization phase with
counter $0$, running exactly $k$ transitions writes, one cell per step while moving
right, the encoded marked blank block over the next $k$ cells, ending in the
initialization phase with counter $k$ and with the work head advanced by $k$ cells; the
input head, output, and remaining tape contents are unchanged.
\end{lemma}

\begin{lemma}[Left extension writes one blank block]
\lean{Turing.FinTM.sweep_grow_left_block}
\label{Turing.FinTM.sweep_grow_left_block}
\leanok
\uses{Turing.Action.apply, Turing.FinTM, Turing.FinTM.SweepAlphabet, Turing.FinTM.blankRow, Turing.FinTM.sweepRevCfg, Turing.FinTM.sweepRevCfg_left_any, Turing.FinTM.sweepSymbol, Turing.FinTM.sweepTM, Turing.FinTM.sweep_generate, Turing.MultiTapeTM.runFrom, Turing.MultiTapeTM.step}
\difficulty{4}
For the single-tape simulator of a machine $M$ with $k$ work tapes and any source state
$q$: from any backward-facing zipper configuration whose control is in the
left-extension phase carrying $q$ with counter $0$, running exactly $k$ transitions
writes the encoded unmarked blank block in reversed cell order while moving left, ending
in the left-extension phase with counter $k$ and the work head moved left by $k$ cells;
the input head and output are unchanged.
\end{lemma}

\begin{definition}[Right guard block]
\lean{Turing.FinTM.rightRow}
\label{Turing.FinTM.rightRow}
\leanok
\uses{Turing.FinTM.SweepCell}
The right guard block of $k$ sweep cells determined by a read table $s$: the cell for
tape $i$ carries index $i$, no payload, a cleared head flag, and a left-neighbour flag
equal to the Boolean recorded for tape $i$ in $s$ — the head flag of the last scanned
block.
\end{definition}

\begin{lemma}[Right extension writes the guard block]
\lean{Turing.FinTM.sweep_grow_right_block}
\label{Turing.FinTM.sweep_grow_right_block}
\leanok
\uses{Turing.Action.apply, Turing.FinTM, Turing.FinTM.SweepAlphabet, Turing.FinTM.rightRow, Turing.FinTM.sweepCfg, Turing.FinTM.sweepCfg_right_any, Turing.FinTM.sweepSymbol, Turing.FinTM.sweepTM, Turing.FinTM.sweep_generate, Turing.MultiTapeTM.runFrom}
\difficulty{4}
For the single-tape simulator of a machine $M$ with $k$ work tapes, a source state $q$,
and a read table $s$: from any forward-facing zipper configuration whose control is in
the right-extension phase carrying $q$ and $s$ with counter $0$, running exactly $k$
transitions writes the guard block determined by $s$ while moving right, ending in the
right-extension phase with counter $k$, still carrying $q$ and $s$, with the work head
advanced by $k$ cells; the input head and output are unchanged.
\end{lemma}

\begin{lemma}[Identity sweep]
\lean{Turing.FinTM.sweepFold_id}
\label{Turing.FinTM.sweepFold_id}
\leanok
\uses{Turing.FinTM.sweepFold}
\difficulty{2}
Folding the identity visit rule — which returns its control state and each visited cell
unchanged — over any list of cells returns the initial state and the unchanged list.
\end{lemma}

\begin{lemma}[Stationary write in a backward zipper]
\lean{Turing.FinTM.sweepRevCfg_write}
\label{Turing.FinTM.sweepRevCfg_write}
\leanok
\uses{Turing.Action.apply, Turing.FinTM.sweepAct, Turing.FinTM.sweepRevCfg, Turing.FinTM.sweepTape, Turing.FinTM.sweepTape_write, Turing.moveInputPos_zero}
\difficulty{3}
In a backward-facing zipper configuration (work head at cell $z$, a list $r$ occupying
cells $z, z-1, \dots$, a list $l$ occupying cells $z+1, z+2, \dots$), applying the sweep
action that writes a symbol $b$ at the current cell without moving, leaves the input
head and output unchanged, and enters a state $q'$ replaces the first entry of $r$ by
$b$ and the control state by $q'$, leaving everything else fixed.
\end{lemma}

\begin{lemma}[Initialization reaches the canonical start]
\lean{Turing.FinTM.sweep_init}
\label{Turing.FinTM.sweep_init}
\leanok
\uses{Turing.Action.apply, Turing.Cfg.workTapeSymbols, Turing.FinTM, Turing.FinTM.SweepAlphabet, Turing.FinTM.SweepState, Turing.FinTM.blankRow, Turing.FinTM.headAt, Turing.FinTM.sweepAct, Turing.FinTM.sweepBoundary, Turing.FinTM.sweepCfg, Turing.FinTM.sweepFold_id, Turing.FinTM.sweepPos, Turing.FinTM.sweepRevCfg, Turing.FinTM.sweepRevCfg_write, Turing.FinTM.sweepStart, Turing.FinTM.sweepSymbol, Turing.FinTM.sweepTM, Turing.FinTM.sweepTape, Turing.FinTM.sweep_init_block, Turing.FinTM.sweep_run_reverse, Turing.FinTM.sweep_turn_left, Turing.FinTM.tapeRow, Turing.FinTM.tapeZone, Turing.MultiTapeTM.initCfg, Turing.MultiTapeTM.runFrom, Turing.MultiTapeTM.runFrom_add, Turing.MultiTapeTM.runFrom_succ_eq_step', Turing.MultiTapeTM.step, Turing.moveInputPos_zero}
\difficulty{5}
Let $M$ be a machine with $k$ work tapes over a finite alphabet $\Gamma$ and let $x$ be
an input word. Running the single-tape simulator for exactly $2k+2$ transitions from its
initial configuration on the symbol-wise encoded input reaches the canonical
between-steps configuration representing the initial configuration of $M$ on $x$, with a
zone of one block at coordinate $0$ and work head at absolute cell $-1$.
\end{lemma}

\begin{lemma}[Stationary phase change, forward form]
\lean{Turing.FinTM.sweepCfg_stay}
\label{Turing.FinTM.sweepCfg_stay}
\leanok
\uses{Turing.Action.apply, Turing.FinTM.sweepCfg, Turing.FinTM.sweepMove}
\difficulty{1}
Applying a stationary phase-change action — no write, no head movement, no output — to a
forward-facing zipper configuration replaces only the control state and leaves the tape,
input head, and output unchanged.
\end{lemma}

\begin{lemma}[Stationary phase change, backward form]
\lean{Turing.FinTM.sweepRevCfg_stay}
\label{Turing.FinTM.sweepRevCfg_stay}
\leanok
\uses{Turing.Action.apply, Turing.FinTM.sweepMove, Turing.FinTM.sweepRevCfg}
\difficulty{1}
Applying a stationary phase-change action — no write, no head movement, no output — to a
backward-facing zipper configuration replaces only the control state and leaves the
tape, input head, and output unchanged.
\end{lemma}

\begin{lemma}[First half of a simulated step]
\lean{Turing.FinTM.sweep_prepare}
\label{Turing.FinTM.sweep_prepare}
\leanok
\uses{Turing.Action.apply, Turing.Cfg, Turing.Cfg.workTapeSymbols, Turing.FinTM, Turing.FinTM.SweepAlphabet, Turing.FinTM.blankRow, Turing.FinTM.headAt, Turing.FinTM.readState, Turing.FinTM.readVisit, Turing.FinTM.read_zone, Turing.FinTM.sweepBoundary, Turing.FinTM.sweepCfg, Turing.FinTM.sweepCfg_stay, Turing.FinTM.sweepFold, Turing.FinTM.sweepPos, Turing.FinTM.sweepRevCfg, Turing.FinTM.sweepStart, Turing.FinTM.sweepSymbol, Turing.FinTM.sweepTM, Turing.FinTM.sweepTape_read, Turing.FinTM.sweep_grow_left_block, Turing.FinTM.sweep_run, Turing.FinTM.sweep_turn_right, Turing.FinTM.tapeRow, Turing.FinTM.tapeRow_blank, Turing.FinTM.tapeZone, Turing.FinTM.tapeZone_length, Turing.MultiTapeTM.runFrom, Turing.MultiTapeTM.runFrom_add, Turing.MultiTapeTM.runFrom_succ_eq_step', Turing.MultiTapeTM.step}
\difficulty{6}
Let $M$ be a machine with $k \ge 1$ work tapes and let $c$ be a configuration of $M$ in
a non-halted state $q$ such that every work-tape head of $c$ is at a coordinate at least
$a$ and every work tape of $c$ is blank at $a-1$. From the canonical between-steps
configuration for $c$ with a zone of $n$ blocks starting at $a$, running the simulator
for exactly $k + 1 + (n+1)k + 1$ transitions extends the zone by a blank block at
$a-1$, performs the forward sweep over all $n+1$ blocks — collecting the scanned symbol
of every source tape into the control and recording left-neighbour head flags in the
cells — and ends in the right-extension phase carrying $q$ and the completed read table
for coordinate $a+n$, with the read-marked encoded blocks enclosed by boundary markers
on the tape and the input head and encoded output preserved.
\end{lemma}

\begin{lemma}[Second half of a simulated step]
\lean{Turing.FinTM.sweep_finish}
\label{Turing.FinTM.sweep_finish}
\leanok
\uses{Turing.Action, Turing.Action.apply, Turing.Cfg, Turing.Cfg.inputSymbol, Turing.Cfg.workTapeSymbols, Turing.FinTM, Turing.FinTM.SweepAlphabet, Turing.FinTM.SweepCell, Turing.FinTM.SweepState, Turing.FinTM.headAt, Turing.FinTM.readState, Turing.FinTM.rightRow, Turing.FinTM.sweepBoundary, Turing.FinTM.sweepCfg, Turing.FinTM.sweepFold, Turing.FinTM.sweepInput_read, Turing.FinTM.sweepPos, Turing.FinTM.sweepPos_move, Turing.FinTM.sweepRevCfg, Turing.FinTM.sweepRevCfg_stay, Turing.FinTM.sweepStart, Turing.FinTM.sweepSymbol, Turing.FinTM.sweepTM, Turing.FinTM.sweepTape_read, Turing.FinTM.sweep_grow_right_block, Turing.FinTM.sweep_run_reverse, Turing.FinTM.sweep_turn_left, Turing.FinTM.tapeRow, Turing.FinTM.tapeZone, Turing.FinTM.tapeZone_append, Turing.FinTM.tapeZone_length, Turing.FinTM.writeVisit, Turing.FinTM.write_zone, Turing.MultiTapeTM.runFrom, Turing.MultiTapeTM.runFrom_add, Turing.MultiTapeTM.runFrom_succ_eq_step', Turing.MultiTapeTM.step, Turing.moveInputPos}
\difficulty{6}
Let $M$ be a machine with $k \ge 1$ work tapes and let $c$ be a configuration with a
state $q$ such that every work-tape head of $c$ is at a coordinate strictly less than
$a+n$ and every work tape of $c$ is blank at $a+n$. From the configuration reached after
the forward sweep — the right-extension phase carrying $q$ and the read table for
coordinate $a+n$, with the read-marked encoded blocks for coordinates
$a-1, \dots, a+n-1$ on the tape between boundary markers — running exactly
$k + 1 + (n+2)k + 1$ transitions writes the right guard block, performs the transition
of $M$ from $q$ on the scanned input symbol and the collected work-tape symbols (moving
the input head and emitting output accordingly), sweeps backward applying that action to
every block, and reaches the canonical between-steps configuration of the resulting
source configuration with a zone of $n+2$ blocks starting at $a-1$.
\end{lemma}

\begin{lemma}[One source step in one bounded burst]
\lean{Turing.FinTM.sweep_step}
\label{Turing.FinTM.sweep_step}
\leanok
\uses{Turing.Cfg, Turing.FinTM, Turing.FinTM.sweepStart, Turing.FinTM.sweepTM, Turing.FinTM.sweepTime, Turing.FinTM.sweep_finish, Turing.FinTM.sweep_prepare, Turing.MultiTapeTM.runFrom, Turing.MultiTapeTM.runFrom_add, Turing.MultiTapeTM.step}
\difficulty{3}
Let $M$ be a machine with $k \ge 1$ work tapes and let $c$ be a configuration in a
non-halted state $q$ whose work-tape heads all lie in the interval $[a, a+n)$ and whose
work tapes are blank at $a-1$ and at $a+n$. Then, from the canonical between-steps
configuration for $c$ with a zone of $n$ blocks starting at $a$, exactly $(2n+5)k + 4$
transitions of the single-tape simulator reach the canonical between-steps configuration
for the configuration obtained by one step of $M$ from $c$, with a zone of $n+2$ blocks
starting at $a-1$ and the absolute work-head position lowered by $k$.
\end{lemma}

\begin{definition}[Exact simulation time]
\lean{Turing.FinTM.sweepTime}
\label{Turing.FinTM.sweepTime}
\leanok
The exact number of simulator transitions elapsed after $t$ source steps, as a function
$\sigma_k : \mathbb{N} \to \mathbb{N}$ defined recursively by $\sigma_k(0) = 2k+2$ and
$\sigma_k(t+1) = \sigma_k(t) + (4t+7)k + 4$.
\end{definition}

\begin{lemma}[Closed form of the simulation time]
\lean{Turing.FinTM.sweepTime_eq}
\label{Turing.FinTM.sweepTime_eq}
\leanok
\uses{Turing.FinTM.sweepTime}
\difficulty{2}
The function $\sigma_k$ defined by $\sigma_k(0) = 2k+2$ and
$\sigma_k(t+1) = \sigma_k(t) + (4t+7)k + 4$ satisfies, for all $k, t \in \mathbb{N}$,
\[
\sigma_k(t) = 2k\,t^2 + (5k+4)\,t + (2k+2).
\]
\end{lemma}

\begin{lemma}[Quadratic bound on the simulation time]
\lean{Turing.FinTM.sweepTime_le}
\label{Turing.FinTM.sweepTime_le}
\leanok
\uses{Turing.FinTM.sweepTime, Turing.FinTM.sweepTime_eq}
\difficulty{4}
The function $\sigma_k$ defined by $\sigma_k(0) = 2k+2$ and
$\sigma_k(t+1) = \sigma_k(t) + (4t+7)k + 4$ satisfies
$\sigma_k(t) \le (9k+6)(t+1)^2$ for all $k, t \in \mathbb{N}$.
\end{lemma}

\begin{lemma}[Simulation up to the first halt]
\lean{Turing.FinTM.sweep_run_to_halt}
\label{Turing.FinTM.sweep_run_to_halt}
\leanok
\uses{Turing.FinTM, Turing.FinTM.SweepAlphabet, Turing.FinTM.source_bounds, Turing.FinTM.sweepStart, Turing.FinTM.sweepTM, Turing.FinTM.sweepTime, Turing.FinTM.sweep_init, Turing.FinTM.sweep_step, Turing.MultiTapeTM.initCfg, Turing.MultiTapeTM.runFrom, Turing.MultiTapeTM.runFrom_add, Turing.MultiTapeTM.runFrom_succ_eq_step'}
\difficulty{5}
Let $M$ be a machine with $k \ge 1$ work tapes over a finite alphabet, let $x$ be an
input word, and suppose the run of $M$ from its initial configuration on $x$ has not
halted at any time $t < \tau$. Then for every $t \le \tau$, running the single-tape
simulator for exactly $\sigma_k(t)$ transitions (the exact simulation time after $t$
source steps) from its initial configuration on the symbol-wise encoded input reaches
the canonical between-steps configuration representing the configuration of $M$ after
$t$ steps, with a zone of $2t+1$ blocks starting at coordinate $-t$ and work head at
absolute cell $-tk - 1$.
\end{lemma}

\begin{theorem}[One work tape suffices]
\lean{Turing.FinTM.one_work_tape}
\label{Turing.FinTM.one_work_tape}
\leanok
\uses{Turing.FinTM, Turing.FinTM.ComputesFunInTime, Turing.FinTM.ComputesFunInTimeVia, Turing.FinTM.ComputesInTime, Turing.FinTM.ComputesInTime.mono, Turing.FinTM.SweepAlphabet, Turing.FinTM.computesInTime_iff, Turing.FinTM.sweepTM, Turing.FinTM.sweepTime, Turing.FinTM.sweepTime_le, Turing.FinTM.sweep_run_to_halt, Turing.FinTM.unusedTapeTM, Turing.FinTM.unusedTape_computes, Turing.MultiTapeTM.initCfg, Turing.MultiTapeTM.runFrom, Turing.MultiTapeTM.runFrom_output_eq_of_halt}
\difficulty{5}
In-model analogue of [AB09, Claim 1.6]. Let $\Gamma$ be a finite alphabet and let $M$ be
a Turing machine over $\Gamma$ (read-only input tape, write-only output tape, any number
of work tapes) that computes a function $f$ on words over $\Gamma$ within time
$T : \mathbb{N} \to \mathbb{N}$, i.e.\ on every input $x$ the run of $M$ has halted
after $T(\abs{x})$ steps with output $f(x)$. Then there exist a finite alphabet
$\Gamma'$, an embedding $e : \Gamma \hookrightarrow \Gamma'$, a machine $M'$ over
$\Gamma'$ with exactly one work tape, and a constant $c \in \mathbb{N}$ such that, on
every symbol-wise encoded input $e(x)$, the run of $M'$ has halted after
$c \cdot (T(\abs{x}) + 1)^2$ steps with output the symbol-wise encoding $e(f(x))$.
\end{theorem}

\begin{theorem}[One work tape and binary alphabet suffice]
\lean{Turing.FinTM.one_work_tape_binary}
\label{Turing.FinTM.one_work_tape_binary}
\leanok
\uses{Turing.FinTM, Turing.FinTM.ComputesFunInTime, Turing.FinTM.ComputesInTime.mono, Turing.FinTM.alphabet_reduction, Turing.FinTM.one_work_tape}
\difficulty{4}
Let $M$ be a Turing machine over the binary alphabet (read-only input tape, write-only
output tape, any number of work tapes) that computes a function $f$ on binary words
within time $T : \mathbb{N} \to \mathbb{N}$, i.e.\ on every input $x$ the run of $M$ has
halted after $T(\abs{x})$ steps with output $f(x)$. Then there exist a machine $M'$ over
the binary alphabet with exactly one work tape and a constant $c \in \mathbb{N}$ such
that $M'$ computes $f$ within time $n \mapsto c \cdot (T(n) + 1)^2$. This combines the
one-work-tape simulation ([AB09, Claim 1.6]) with alphabet reduction
([AB09, Claim 1.5]).
\end{theorem}
